\section{Introduction}

Professional music production workflows treat MIDI as an editable artifact: notes, velocities, instruments, and tempo are parameters a producer adjusts iteratively. Recent large language models can discuss music theory, transcribe audio, and generate symbolic scores, but there is no standard benchmark that asks a model to \emph{edit} an existing MIDI file under a natural-language instruction while checking both requested changes and explicitly declared preservation constraints.

Existing work clusters into three nearby problems, none of which matches this task:
\begin{enumerate}[leftmargin=*]
\item \textbf{Symbolic reasoning benchmarks} such as WildScore and MusicTheoryBench evaluate multiple-choice understanding over score images or theory questions, not executable edits.
\item \textbf{Audio editing systems} such as InstructME and MuseCPBench evaluate remixing and context preservation in waveform space; RefinPaint~\cite{ramoneda2024refinpaint} performs symbolic proofreading but is not conditioned on free-text edits.
\item \textbf{Zero-shot drum editing} (Not that Groove) demonstrates instruction-driven symbolic edits with unit tests, but covers one-bar drum grooves without a preservation axis.
\end{enumerate}

MIDI-Instruct fills the gap: \emph{MIDI in, instruction in, MIDI out}, with dual metrics that mirror Music Context Preservation (MCP) but operate on symbolic features rather than rendered audio.

\paragraph{Contributions.}
\begin{enumerate}[leftmargin=*]
\item \textbf{MIDI-Instruct}, a JSONL benchmark format pairing MIDI, instructions, edit masks, and predicate specifications (Section~3).
\item A \textbf{dual-metric grader} reporting \edit{}, \preserve{}, joint pass rate, over-edit rate, and optional gold note-F1 (Section~4).
\item \textbf{MIDI-Reason}, a structured plan track scored independently from the output file (Section~3.3).
\item An open \textbf{generation and evaluation harness} producing 300 unique-gold pilot items from programmatic seeds, with reference and Llama-3.2-1B baselines (Sections~5--7).
\end{enumerate}

Our first model evaluation reaches 0.44 joint pass rate while exact plan match is only 0.20. The large gap across operator families---perfect performance on muting and tempo scaling but no joint passes on velocity or program changes---demonstrates why aggregate output similarity alone is insufficient. Later work---text-to-MIDI generation, tool-augmented agents, and self-refining harnesses---can reuse the same item schema; only the runner changes.
